%==============================================================================
% Sjabloon onderzoeksvoorstel bachproef
%==============================================================================
% Gebaseerd op document class `hogent-article'
% zie <https://github.com/HoGentTIN/latex-hogent-article>

% Voor een voorstel in het Engels: voeg de documentclass-optie [english] toe.
% Let op: kan enkel na toestemming van de bachelorproefcoördinator!
\documentclass{hogent-article}

% Invoegen bibliografiebestand
\addbibresource{voorstel.bib}

% Informatie over de opleiding, het vak en soort opdracht
\studyprogramme{Professionele bachelor toegepaste informatica}
\course{Bachelorproef}
\assignmenttype{Onderzoeksvoorstel}
% Voor een voorstel in het Engels, haal de volgende 3 regels uit commentaar
% \studyprogramme{Bachelor of applied information technology}
% \course{Bachelor thesis}
% \assignmenttype{Research proposal}

\academicyear{2025-2026}

% TODO: Werktitel
\title{Ontwerp en realisatie van een datagedreven scoutingomgeving voor de BNXT League}

\author{Seppe Baert}
\email{Seppe.Baert@student.hogent.be}


% TODO: Geef de co-promotor op
% \supervisor[Co-promotor]{S. Beekman (Synalco, \href{mailto:sigrid.beekman@synalco.be}{sigrid.beekman@synalco.be})}


% Binnen welke specialisatierichting uit 3TI situeert dit onderzoek zich?
% Kies uit deze lijst:
%
% - Mobile \& Enterprise development
% - AI \& Data Engineering
% - Functional \& Business Analysis
% - System \& Network Administrator
% - Mainframe Expert
% - Als het onderzoek niet past binnen een van deze domeinen specifieer je deze
%   zelf
%
\specialisation{AI \& Data Engineering}
\keywords{BNXT League, Data analyse, Scouting}

\begin{document}
\begin{abstract}
Deze bachelorproef richt zich op het ontwerpen en realiseren van een datagedreven scoutingomgeving voor de BNXT League, een competitie waar clubs wel over wedstrijd- en spelersstatistieken beschikken maar waar rekruteringsbeslissingen vaak nog gebaseerd worden op intuïtie, losse bronnen en persoonlijke netwerken. De probleemstelling vertrekt vanuit deze kloof tussen beschikbare data en het ontbreken van een geïntegreerd, herhaalbaar beslissingskader voor spelersselectie, ondanks de aanwezigheid van relevante prestatie-indicatoren in bestaande basketbalstatistieken. De centrale onderzoeksvraag luidt: Hoe kan een datagedreven scoutingomgeving voor de BNXT League worden ontworpen en gerealiseerd zodat scouts spelers efficiënt kunnen analyseren en vergelijken op basis van relevante prestatie-indicatoren? 
\\\\
Het concrete doel is een proof‑of‑concept op te leveren met een centrale datalaag en een Power BI dashboard waarmee BNXT spelers gefilterd, gerangschikt en onderling vergeleken kunnen worden, aangevuld met documentatie van datamodel en indicatoren. Methodologisch wordt eerst via een literatuurstudie bepaald welke statistieken en besluitvormingskaders relevant zijn, gevolgd door het verzamelen, opschonen en modelleren van BNXT data en de implementatie van een Power BI dashboard met spelersoverzichten, detailpagina’s en vergelijkingsschermen. De omgeving wordt geëvalueerd via technische tests en een beperkte gebruiksbeoordeling aan de hand van vooraf gedefinieerde use cases. Verwacht wordt dat deze scoutingomgeving een deel van de huidige, eerder intuïtieve besluitvorming vervangt door gestructureerde en herhaalbare analyses, zodat scouts en sportief managers beter onderbouwde beslissingen kunnen nemen binnen de BNXT context.
\end{abstract}

\tableofcontents

% De hoofdtekst van het voorstel zit in een apart bestand, zodat het makkelijk
% kan opgenomen worden in de bijlagen van de bachelorproef zelf.
%---------- Inleiding ---------------------------------------------------------

% TODO: Is dit voorstel gebaseerd op een paper van Research Methods die je
% vorig jaar hebt ingediend? Heb je daarbij eventueel samengewerkt met een
% andere student?
% Zo ja, haal dan de tekst hieronder uit commentaar en pas aan.

%\paragraph{Opmerking}

% Dit voorstel is gebaseerd op het onderzoeksvoorstel dat werd geschreven in het
% kader van het vak Research Methods dat ik (vorig/dit) academiejaar heb
% uitgewerkt (met medesturent VOORNAAM NAAM als mede-auteur).
% 

\section{Inleiding}%
\label{sec:inleiding}

Basketbal en in het bijzonder de National Basketball Association (NBA), is een ondertussen sterk gedigitaliseerde sport waarin alle gebeurtenissen worden omgezet in data: punten, rebounds, assists, speeltijd, enzovoort. Toch blijft het accuraat voorspellen van spelersprestaties een grote uitdaging omdat er zoveel verschillende factoren meespelen zoals onder andere vorm, spelpositie, tegenstanders, blessures en meer. Doordat er talrijke historische spelersdata is, is het hierdoor wel mogelijk om patronen te ontdekken in deze data. Wat momenteel het probleem vormt is dat eenvoudige modellen deze data niet benuttigen. Maar door de enorme opkomst van Artificiële Intelligentie (AI) de laatste jaren, zijn er al veelbelovende resultaten geboekt. Patronen die voor menselijke analisten moeilijk zichtbaar kunnen zijn.
\\
\\
Dit onderzoek richt zich op het gebruiken en vergelijken van verschillende machine learning modellen om speler statistieken te voorspellen. Dit gebeurt aan de hand van historische wedstrijddata. Het gaat hierbij om dingen zoals het aantal punten, rebounds of assists dat een speler per wedstrijd behaald. De doelgroep van dit onderzoek zijn vooral data analisten binnen sport analytische bedrijven of NBA teams die deze data willen gebruiken om beslissingen te maken eventueel rond scouting of speler evaluatie.
\\
\\
De probleemstelling van dit onderzoek: Hoe presteren verschillende machine learning modellen bij het voorspellen van NBA speler statistieken op basis van historische wedstrijddata en welk model levert daarbij de meest nauwkeurige voorspellingen op? Inmiddels zijn er verschillende modellen die wedstrijduitslagen kunnen voorspellen, maar er is minder helderheid over hoe spelers hun prestaties zullen vertonen. Het is daarom boeiend om te kijken naar de kwaliteitsverschillen tussen deze modellen.
\\
\\
Dit onderzoek heeft als doel om te bepalen welke machine learning-modellen het beste geschikt zijn voor het voorspellen van speler statistieken, gebaseerd op een literatuurstudie en experimentele vergelijkingen. We gaan verschillende modellen trainen en evalueren op dezelfde historische dataset met maatstaven zoals Mean Absolute Error (MAE) en Root Mean Squared Error (RMSE).

%---------- Stand van zaken ---------------------------------------------------

\section{Literatuurstudie}%
\label{sec:literatuurstudie}

Hier beschrijf je de \emph{state-of-the-art} rondom je gekozen onderzoeksdomein, d.w.z.\ een inleidende, doorlopende tekst over het onderzoeksdomein van je bachelorproef. Je steunt daarbij heel sterk op de professionele \emph{vakliteratuur}, en niet zozeer op populariserende teksten voor een breed publiek. Wat is de huidige stand van zaken in dit domein, en wat zijn nog eventuele open vragen (die misschien de aanleiding waren tot je onderzoeksvraag!)?

Je mag de titel van deze sectie ook aanpassen (literatuurstudie, stand van zaken, enz.). Zijn er al gelijkaardige onderzoeken gevoerd? Wat concluderen ze? Wat is het verschil met jouw onderzoek?

Verwijs bij elke introductie van een term of bewering over het domein naar de vakliteratuur, bijvoorbeeld~\autocite{Hykes2013}! Denk zeker goed na welke werken je refereert en waarom.

Draag zorg voor correcte literatuurverwijzingen! Een bronvermelding hoort thuis \emph{binnen} de zin waar je je op die bron baseert, dus niet er buiten! Maak meteen een verwijzing als je gebruik maakt van een bron. Doe dit dus \emph{niet} aan het einde van een lange paragraaf. Baseer nooit teveel aansluitende tekst op eenzelfde bron.

Als je informatie over bronnen verzamelt in JabRef, zorg er dan voor dat alle nodige info aanwezig is om de bron terug te vinden (zoals uitvoerig besproken in de lessen Research Methods).

% Voor literatuurverwijzingen zijn er twee belangrijke commando's:
% \autocite{KEY} => (Auteur, jaartal) Gebruik dit als de naam van de auteur
%   geen onderdeel is van de zin.
% \textcite{KEY} => Auteur (jaartal)  Gebruik dit als de auteursnaam wel een
%   functie heeft in de zin (bv. ``Uit onderzoek door Doll & Hill (1954) bleek
%   ...'')

Je mag deze sectie nog verder onderverdelen in subsecties als dit de structuur van de tekst kan verduidelijken.

%---------- Methodologie ------------------------------------------------------
\section{Methodologie}%
\label{sec:methodologie}

Hier beschrijf je hoe je van plan bent het onderzoek te voeren. Welke onderzoekstechniek ga je toepassen om elk van je onderzoeksvragen te beantwoorden? Gebruik je hiervoor literatuurstudie, interviews met belanghebbenden (bv.~voor requirements-analyse), experimenten, simulaties, vergelijkende studie, risico-analyse, PoC, \ldots?

Valt je onderwerp onder één van de typische soorten bachelorproeven die besproken zijn in de lessen Research Methods (bv.\ vergelijkende studie of risico-analyse)? Zorg er dan ook voor dat we duidelijk de verschillende stappen terug vinden die we verwachten in dit soort onderzoek!

Vermijd onderzoekstechnieken die geen objectieve, meetbare resultaten kunnen opleveren. Enquêtes, bijvoorbeeld, zijn voor een bachelorproef informatica meestal \textbf{niet geschikt}. De antwoorden zijn eerder meningen dan feiten en in de praktijk blijkt het ook bijzonder moeilijk om voldoende respondenten te vinden. Studenten die een enquête willen voeren, hebben meestal ook geen goede definitie van de populatie, waardoor ook niet kan aangetoond worden dat eventuele resultaten representatief zijn.

Uit dit onderdeel moet duidelijk naar voor komen dat je bachelorproef ook technisch voldoen\-de diepgang zal bevatten. Het zou niet kloppen als een bachelorproef informatica ook door bv.\ een student marketing zou kunnen uitgevoerd worden.

Je beschrijft ook al welke tools (hardware, software, diensten, \ldots) je denkt hiervoor te gebruiken of te ontwikkelen.

Probeer ook een tijdschatting te maken. Hoe lang zal je met elke fase van je onderzoek bezig zijn en wat zijn de concrete \emph{deliverables} in elke fase?

%---------- Verwachte resultaten ----------------------------------------------
\section{Verwacht resultaat, conclusie}%
\label{sec:verwachte_resultaten}

Hier beschrijf je welke resultaten je verwacht. Als je metingen en simulaties uitvoert, kan je hier al mock-ups maken van de grafieken samen met de verwachte conclusies. Benoem zeker al je assen en de onderdelen van de grafiek die je gaat gebruiken. Dit zorgt ervoor dat je concreet weet welk soort data je moet verzamelen en hoe je die moet meten.

Wat heeft de doelgroep van je onderzoek aan het resultaat? Op welke manier zorgt jouw bachelorproef voor een meerwaarde?

Hier beschrijf je wat je verwacht uit je onderzoek, met de motivatie waarom. Het is \textbf{niet} erg indien uit je onderzoek andere resultaten en conclusies vloeien dan dat je hier beschrijft: het is dan juist interessant om te onderzoeken waarom jouw hypothesen niet overeenkomen met de resultaten.



\printbibliography[heading=bibintoc]

\end{document}